\documentclass[10pt,aspectratio=169]{beamer}

% --- Configuración de Idioma, Codificación e Imágenes ---
\usepackage[spanish]{babel}
\usepackage[utf8]{inputenc}
\usepackage{tikz}
\usepackage{graphicx} % Para la carga de gráficos

% --- Tipografía "Tech Performance" ---
\usefonttheme{serif} % Títulos en Serif (Ingeniería Académica)
\usepackage{inconsolata} % Fuente Monospace para subtítulos y código

% --- Configuración de Colores (Modo Oscuro / Terminal) ---
\definecolor{TechBg}{HTML}{0D1117}       % Fondo Negro Profundo (Estilo IDE)
\definecolor{TechRed}{HTML}{FF3333}      % Rojo Código / Acento
\definecolor{TechWhite}{HTML}{F8F9FA}    % Texto Principal
\definecolor{TechGray}{HTML}{8B949E}     % Texto Secundario/Comentarios

% Aplicación de colores globales
\setbeamercolor{background canvas}{bg=TechBg}
\setbeamercolor{normal text}{fg=TechWhite}
\setbeamercolor{structure}{fg=TechRed}
\setbeamercolor{title}{fg=TechWhite}
\setbeamercolor{subtitle}{fg=TechGray}
\setbeamercolor{author}{fg=TechWhite}
\setbeamercolor{date}{fg=TechGray}

% --- Estilo de Elementos y Listas ---
\setbeamertemplate{navigation symbols}{}
\setbeamercolor{item}{fg=TechRed}
\setbeamertemplate{itemize item}{\small$\blacksquare$}
\setbeamertemplate{itemize subitem}{\scriptsize$\rightarrow$}

% --- Plantilla de Títulos de Diapositiva ---
\setbeamertemplate{frametitle}{
    \vspace{0.4cm}
    \begin{minipage}{\textwidth}
        {\usebeamerfont{frametitle}\bfseries\insertframetitle}
        \vskip2pt
        \tikz\draw[TechRed, thick] (0,0) -- (\textwidth,0);
    \end{minipage}
}

% --- Datos de la Presentación ---
\title{\Huge Divide y Vencerás}
\subtitle{\texttt{[Análisis Experimental] Ordenamiento, Búsqueda y Selección}}
\author{A. Barbosa \and E. Velásquez \and D. Oyarzo}
\date{\texttt{Mayo 2026 // C\_CORE\_PERFORMANCE}}

\begin{document}

% --- DIAPOSITIVA DE PORTADA ---
\begin{frame}[plain]
    \vspace{1.5cm}
    \begin{columns}
        \begin{column}{0.65\textwidth}
            {\usebeamercolor[fg]{title}\usebeamerfont{title}\inserttitle}\vspace{0.2cm}
            
            {\usebeamercolor[fg]{subtitle}\fontfamily{zi4}\selectfont\insertsubtitle}\vspace{1cm}
            
            \tikz\fill[TechRed] (0,0) rectangle (2cm, 0.1cm);
            \vspace{0.5cm}
            
            {\small \textbf{Investigadores:} \insertauthor}
        \end{column}
        
        \begin{column}{0.3\textwidth}
            \flushright
            \small
            \color{TechGray}
            \texttt{SYS: GESTIÓN\_DEPORTISTAS}\\
            \texttt{LANG: C NATIVO}\\
            \texttt{\insertdate}
        \end{column}
    \end{columns}
\end{frame}

%----------------------------------------
\begin{frame}{Objetivo del Proyecto}
\begin{itemize}
\item Analizar empírica y teóricamente los algoritmos basados en \textbf{Divide y Vencerás}.
\item Aplicar los modelos sobre un sistema robusto y nativo de gestión en lenguaje \texttt{C}.
\item Validar el impacto de las heurísticas a través de tres áreas core:
\begin{itemize}
\item \texttt{ORDENAMIENTO} $\rightarrow$ Merge Sort (Clásico/Threshold), Quick Sort (Lomuto).
\item \texttt{BÚSQUEDA} $\rightarrow$ Binaria, Exponencial, Interpolación.
\item \texttt{SELECCIÓN} $\rightarrow$ Quick Select (búsqueda del $k$-ésimo elemento).
\end{itemize}
\end{itemize}

\vspace{0.6cm}
\textbf{\color{TechRed}\texttt{>>}} \textbf{Meta Asintótica:} Evidenciar rigurosamente el salto temporal de $\mathcal{O}(n^2)$ a $\mathcal{O}(n \log n)$.
\end{frame}

%----------------------------------------
\begin{frame}{Algoritmos de Ordenamiento: Base Teórica}
\begin{columns}[t]
    \begin{column}{0.48\textwidth}
        \textbf{\color{TechRed} \texttt{01\_}} \textbf{Merge Sort}
        \vskip5pt
        {\footnotesize 
        \textbf{Clásico:} División recursiva sistemática hasta $n=1$. Garantiza $\mathcal{O}(n \log n)$ en todos los escenarios, a costo de uso extra de memoria $O(n)$.\\[1ex]
        \textbf{Optimizado:} Delega subarreglos pequeños a un \texttt{Insertion Sort} controlando el límite mediante un \textit{threshold}. Disminuye radicalmente el \textit{overhead} del stack recursivo.
        }
    \end{column}
    \begin{column}{0.48\textwidth}
        \textbf{\color{TechRed} \texttt{02\_}} \textbf{Quick Sort}
        \vskip5pt
        {\footnotesize 
        \textbf{Mecánica:} Basado puramente en el esquema de partición de Lomuto.\\[1ex]
        \textbf{Impacto del Pivote:} Se analizaron múltiples heurísticas (Aleatorio, Mediana, etc.) para evitar el peor caso $\mathcal{O}(n^2)$ que ocurre al procesar arreglos previamente ordenados.
        }
    \end{column}
\end{columns}
\end{frame}

%----------------------------------------
\begin{frame}{Merge Sort: Clásico vs Optimizado}
\begin{columns}
    \begin{column}{0.48\textwidth}
        \centering
        % Reemplazar "figura1.png" por el nombre exacto de tu imagen compilada
        \includegraphics[width=\textwidth,height=0.55\textheight,keepaspectratio]{figura1.png}\\
        {\scriptsize \textbf{Figura 1:} Merge Sort (Todos los casos)}
    \end{column}
    \begin{column}{0.48\textwidth}
        \centering
        % Reemplazar "figura5.png" por el nombre exacto de tu imagen compilada
        \includegraphics[width=\textwidth,height=0.55\textheight,keepaspectratio]{figura5.png}\\
        {\scriptsize \textbf{Figura 5:} Merge Sort Optimizado}
    \end{column}
\end{columns}
\end{frame}

%----------------------------------------
\begin{frame}{Optimización de Umbral (\textit{Threshold}) en Merge Sort}
\begin{columns}
    \begin{column}{0.31\textwidth}
        \centering
        \includegraphics[width=\textwidth]{figura3.png}\\
        {\scriptsize \textbf{Fig 3:} Mejor caso}
    \end{column}
    \begin{column}{0.31\textwidth}
        \centering
        \includegraphics[width=\textwidth]{figura2.png}\\
        {\scriptsize \textbf{Fig 2:} Caso Promedio}
    \end{column}
    \begin{column}{0.31\textwidth}
        \centering
        \includegraphics[width=\textwidth]{figura4.png}\\
        {\scriptsize \textbf{Fig 4:} Peor caso}
    \end{column}
\end{columns}

\vspace{0.4cm}
\centering{\footnotesize\textbf{\texttt{[ANALÍTICA]}} La selección paramétrica del tamaño de combinación (threshold) moldea determinantemente el uso eficiente de CPU en conjuntos de datos pequeños.}
\end{frame}

%----------------------------------------
\begin{frame}{Quick Sort: Comparación y Degradación}
\begin{columns}
    \begin{column}{0.48\textwidth}
        \centering
        \includegraphics[width=\textwidth,height=0.55\textheight,keepaspectratio]{figura7.png}\\
        {\scriptsize \textbf{Figura 7:} Quick Sort (Caso Promedio)}
    \end{column}
    \begin{column}{0.48\textwidth}
        \centering
        \includegraphics[width=\textwidth,height=0.55\textheight,keepaspectratio]{figura6.png}\\
        {\scriptsize \textbf{Figura 6:} Quick Sort (Peor Caso)}
    \end{column}
\end{columns}
\end{frame}

%----------------------------------------
\begin{frame}{Benchmarking Global de Ordenamiento}
\begin{columns}
    \begin{column}{0.48\textwidth}
        \centering
        \includegraphics[width=\textwidth,height=0.55\textheight,keepaspectratio]{figura9.png}\\
        {\scriptsize \textbf{Figura 9:} Comparativa (Caso Promedio)}
    \end{column}
    \begin{column}{0.48\textwidth}
        \centering
        \includegraphics[width=\textwidth,height=0.55\textheight,keepaspectratio]{figura10.png}\\
        {\scriptsize \textbf{Figura 10:} Comparativa (Peor Caso)}
    \end{column}
\end{columns}
\end{frame}

%----------------------------------------
\begin{frame}{Algoritmos de Búsqueda y Selección}
\begin{itemize}
\item \textbf{Búsqueda Binaria:} Salto de $\mathcal{O}(n)$ a $\mathcal{O}(\log n)$. Evaluada mediante recursión nativa y técnica de delimitación de rangos directos.
\item \textbf{Búsqueda Exponencial:} Expansión del rango de búsqueda mediante potencias de $2^i$, finalizando en una búsqueda binaria local.
\item \textbf{Búsqueda por Interpolación:} Estimación directa de la posición evaluando valores de los extremos. Óptima en distribuciones uniformes ($\mathcal{O}(\log \log n)$).
\end{itemize}

\vspace{0.4cm}
\textbf{\color{TechRed}Extracción de Datos:}
\begin{itemize}
\item \textbf{Quick Select:} Recuperación del $k$-ésimo elemento de un conjunto desordenado descartando recursivamente la mitad inútil. $\mathcal{O}(n)$ en el caso promedio.
\end{itemize}
\end{frame}

%----------------------------------------
\begin{frame}{Rendimiento en Algoritmos de Búsqueda}
\begin{columns}
    \begin{column}{0.48\textwidth}
        \centering
        \includegraphics[width=\textwidth,height=0.55\textheight,keepaspectratio]{figura11.png}\\
        {\scriptsize \textbf{Figura 11:} Búsqueda (Caso Promedio)}
    \end{column}
    \begin{column}{0.48\textwidth}
        \centering
        \includegraphics[width=\textwidth,height=0.55\textheight,keepaspectratio]{figura12.png}\\
        {\scriptsize \textbf{Figura 12:} Búsqueda (Peor Caso)}
    \end{column}
\end{columns}
\end{frame}

%----------------------------------------
\begin{frame}{Comportamiento de Quick Select}
\begin{columns}
    \begin{column}{0.5\textwidth}
        \centering
        \includegraphics[width=0.9\textwidth]{figura13.png}\\
        {\scriptsize \textbf{Figura 13:} Rendimiento en Extracción (Mejor y Peor Caso)}
    \end{column}
    \begin{column}{0.5\textwidth}
        \textbf{\color{TechRed}Análisis Operativo:}
        \begin{itemize}
        \item Se comporta idéntico a Quick Sort, pero al desechar particiones logra una complejidad lineal $\mathcal{O}(n)$.
        \item Ideal en escenarios de métricas deportivas donde es necesario extraer ránkings (el $top-k$) sin ordenar todo el set masivo de datos.
        \end{itemize}
    \end{column}
\end{columns}
\end{frame}

%----------------------------------------
\begin{frame}{Conclusión}
\centering
\begin{tikzpicture}[node distance=2.5cm, auto]
    \tikzstyle{block} = [rectangle, draw=TechRed, fill=TechBg, text=TechWhite, width=2.8cm, text centered, minimum height=1cm, ultra thick]
    \tikzstyle{line} = [draw, -latex, TechWhite, thick]
    
    \node [block] (teoria) {\footnotesize \textbf{Base Teórica}};
    \node [block, right of=teoria, xshift=1.5cm] (impl) {\footnotesize \textbf{Optimización C}};
    \node [block, right of=impl, xshift=1.5cm] (bench) {\footnotesize \textbf{Benchmarking}};
    
    \path [line] (teoria) -- (impl);
    \path [line] (impl) -- (bench);
\end{tikzpicture}

\vspace{0.8cm}
\begin{itemize}
\item La gestión directa de memoria (\texttt{C nativo}) revela el peso real del \textit{overhead} recursivo que la cota asintótica grande $\mathcal{O}$ normalmente ignora.
\item El paradigma \textbf{Divide y Vencerás} requiere afinación manual (como el esquema híbrido Merge/Insertion o la heurística del pivote) para no colapsar el rendimiento en entornos de producción.
\end{itemize}

\vspace{0.6cm}
\centering{\footnotesize\color{TechGray}\texttt{[FIN DEL ANÁLISIS // MEMORY\_FREED]}}
\end{frame}

\end{document}